\documentclass[12pt,a4paper]{article}
\usepackage[utf8]{inputenc}
\usepackage[T2A]{fontenc}
\usepackage[russian]{babel}
\usepackage{amsmath,amssymb,amsfonts,mathtools}
\usepackage{geometry}
\geometry{top=2cm, bottom=2cm, left=2.5cm, right=2cm}
\usepackage{enumitem}
\usepackage{hyperref}
\usepackage{amsthm}
\theoremstyle{definition}
\newtheorem{definition}{Определение}[section]
\usepackage{mathrsfs}
\hypersetup{colorlinks=true,linkcolor=blue,urlcolor=blue}


\begin{document}

\title{Домашняя работа по теории вероятностей\\ Вариант №5}
\author{Студент: Виноградов Илья \quad Группа: СКБ242}
\maketitle

\section{Распределения}
\begin{itemize}
    \item \textbf{Дискретное распределение $\mathscr{L}_d$ (№5):} Геометрическое распределение $\xi \sim \text{Geom}(\theta)$
    \[
    P(\xi = x) = \theta (1-\theta)^{x-1}, \quad x \in \mathbb{N}, \quad 0 < \theta < 1
    \]
    \item \textbf{Непрерывное распределение $\mathscr{L}_c$ (№5):} Треугольное распределение (Симпсона) $\xi \sim \text{Triang}(0,1,\theta)$, $\theta \in (0,1)$
    \[
    f_\xi(x) = 
    \begin{cases}
        \dfrac{2x}{\theta}, & 0 \le x \le \theta,\\[8pt]
        \dfrac{2(1-x)}{1-\theta}, & \theta < x \le 1,\\[8pt]
        0, & \text{иначе}
    \end{cases}
    \]
\end{itemize}

\renewcommand{\thesection}{1.\arabic{section}}
\setcounter{section}{0}
\section{Основные понятия и определения}

\begin{definition}[Функция распределения]
Функция $F(x)$ называется \textbf{функцией распределения}, если выполнены следующие условия:
\begin{enumerate}[label=\arabic*.]
    \item $F(x)$ не убывает.
    \item $F(x)$ непрерывна слева в каждой точке $x \in \mathbb{R}$.
    \item $F(-\infty) = 0$, $F(+\infty) = 1$.
\end{enumerate}
\end{definition}

\begin{definition}[Квантиль распредления уровня $\gamma$]
Число $t_{\gamma}$ называется \textbf{квантилем распределения уровня $\gamma$ для с.в $\xi$}, если {$F_\xi (t_\gamma) = \gamma$} и $t_{\gamma}$:
\begin{enumerate}[label=\arabic*.]
    \item $F_{\xi}(t_{\gamma}) \le \gamma$
    \item $P(\xi > t_{\gamma}) \ge 1 - \gamma$
\end{enumerate}
\end{definition}

\begin{definition}[Медиана распределения]
\textbf{Медианой распределения} называется квантиль распределения уровня $\frac{1}{2}$
\end{definition}

\begin{definition}[Математическое ожидание]
    Математическое ожидание случайной величины введем в три этапа
    \begin{enumerate}[label=\arabic*.]
        \item \textbf{Для простой случайной величины}
        Математическим ожиданием простой случайной величины $\xi = \sum_{i=1}^{n} x_i I(A_i)$ будем обозначать $M\xi$ и называть число, равное
        \[
        M\xi = \int_{\Omega} \xi(\omega) \, P(d\omega) := \sum_{i=1}^{n} x_i P(A_i)
        \]
        
        \item \textbf{Для неотрицательной случайной величины}
        Пусть $\xi \ge 0$ — произвольная неотрицательная случайная величина, определённая на $(\Omega, \mathcal{F}, P)$
        Пусть $\xi_n$ — последовательность простых случайных величин такая, что для произвольного $\omega \in \Omega$ последовательность $\xi_n(\omega)$ монотонно сходится к $\xi(\omega)$:
        \[
        \xi_n(\omega) \uparrow \xi(\omega), \quad n \to \infty
        \]
        Тогда \textbf{математическим ожиданием} случайной величины $\xi$ называется следующий предел:
        \[
        M\xi = \lim_{n \to \infty} M\xi_n
        \]
        
        \item \textbf{Для произвольной случайной величины}
        Пусть $\xi$ — произвольная случайная величина. Её можно представить в виде суммы неотрицательных случайных величин:
        \[
        \xi(\omega) = \xi \cdot I(\xi \ge 0) + \xi \cdot I(\xi < 0)
        \]
        Определим следующие случайные величины:
        \[
        \xi^+ = \xi \cdot I(\xi \ge 0), \qquad \xi^- = -\xi \cdot I(\xi < 0)
        \]
        Тогда $\xi^+$ и $\xi^-$ — неотрицательные случайные величины, для которых определено математическое ожидание (интеграл Лебега):
        \[
        I^+ = M\xi^+, \qquad I^- = M\xi^-
        \]
        
        Если $I^+, I^- < \infty$, то по определению математическое ожидание (интеграл Лебега) случайной величины $\xi$ равно
        \[
        M\xi = I^+ - I^-
        \]
        
        Если хотя бы один из интегралов $I^+$ или $I^-$ расходится, то говорят, что математическое ожидание не определено
    \end{enumerate}
\end{definition}

\begin{definition}[Математическое ожидание функции от случайной величины]
Пусть $g(x)$ --- измеримая функция на $\mathbb{R}$. Величина $Mg(\xi)$ называется \textbf{математическим ожиданием} функции $g$ от случайной величины $\xi$
\end{definition}

\begin{definition}[Момент $n$-го порядка]
Пусть $g(x) = x^n$. Величина $M\xi^n$ называется \textbf{моментом $n$-го порядка}
\begin{itemize}
    \item Для дискретных случайных величин:
    \[
    M\xi^n = \sum_{k} x_k^n \cdot P(\xi = x_k)
    \]
    \item Для абсолютно непрерывных случайных величин:
    \[
    M\xi^n = \int_{\mathbb{R}} x^n p_\xi(x)\,dx
    \]
\end{itemize}
\end{definition}

\begin{definition}[Центральный момент $n$-го порядка]
Пусть $g(x) = (x - \mathbb{E}\xi)^n$. Величина $M(\xi - M\xi)^n$ называется \textbf{центральным моментом $n$-го порядка}
\begin{itemize}
    \item Для дискретной случайной величины:
    \[
    M(\xi - M\xi)^n = \sum_{k} (x_k - M\xi)^n \cdot P(\xi = x_k)
    \]
    \item Для абсолютно непрерывной случайной величины:
    \[
    M(\xi - M\xi)^n = \int_{\mathbb{R}} (x - M\xi)^n \cdot p_\xi(x)\,dx
    \]
\end{itemize}
\end{definition}

\begin{definition}[Дисперсия]
Центральный момент второго порядка называется \textbf{дисперсией} и обозначается $D\xi$ или $\sigma^2(\xi)$:
\[
D\xi = \sigma^2(\xi) = M(\xi - M\xi)^2
\]
\end{definition}

\begin{definition}[Факториальный момент $k$-го порядка]
Величины
\[
M\nu_{[k]} = M\bigl[\nu(\nu-1)\dots(\nu-k+1)\bigr]
\]
называются \textbf{факториальными моментами} целочисленной случайной величины $\nu$
\end{definition}

\begin{definition}[Мода распределения]
\textbf{Модой распределения} называется такое значение случайной величины, в котором достигается максимум вероятности (или плотности):
\begin{itemize}
    \item для абсолютно непрерывной случайной величины: точка максимума плотности $f_\xi(x)$;
    \item для дискретной случайной величины: значение $x_i$, для которого $P(\xi = x_i)$ максимальна
\end{itemize}
\end{definition}

\begin{definition}[Производящая функция]
\textbf{Производящей функцией} $f_\nu(s)$ распределения случайной величины $\nu$, принимающей целые неотрицательные значения, называется степенной ряд (для $|s| < 1$):
\[
f_\nu(s) = Ms^\nu = \sum_{r=0}^{\infty} P(\nu = r)\, s^r
\]
\end{definition}

\begin{definition}[Характеристическая функция]
\textbf{Характеристической функцией} случайной величины $\xi$, принимающей действительные значения, называется функция
\[
f(t) = Me^{it\xi}, \quad t \in \mathbb{R}.
\]
В частности:
\begin{itemize}
    \item для случайной величины $\xi$ с дискретным распределением:
    \[
    f(t) = \sum_{k} e^{it x_k} P(\xi = x_k);
    \]
    \item для случайной величины $\xi$, распределение которой имеет плотность $p(x)$:
    \[
    f(t) = \int_{\mathbb{R}} e^{itx} p(x)\,dx;
    \]
    \item в общем случае для случайной величины $\xi$ с функцией распределения $F(x) = P(\xi \le x)$:
    \[
    f(t) = \int_{\Omega} e^{it\xi(\omega)} P(d\omega) = \int_{\mathbb{R}} e^{itx} dF(x)
    \]
\end{itemize}
\end{definition}

\section{Примеры и взаимосвязи}

\subsection{Геометрическое распределение $\mathscr{L}_d=\mathrm{Geom}(\theta)$}

Пусть случайная величина $\xi$ имеет геометрическое распределение с параметром
$\theta \in (0,1)$:
\[
P(\xi = k) = (1-\theta)^{k-1}\theta,
\qquad k=1,2,\dots
\]

\textbf{Эксперимент}
Пользователь вводит пароль в систему. Вероятность того, что пароль введён правильно с одной попытки, равна $\theta$. Попытки независимы

Случайная величина $\xi$ - номер попытки, на которой пользователь впервые ввёл пароль правильно

\textbf{Обоснование}
Событие $\{\xi=k\}$ означает:
\begin{itemize}
    \item первые $k-1$ попыток были неверными;
    \item на $k$-й попытке пароль введён правильно
\end{itemize}

Так как попытки независимы,
\[
P(\xi=k)
=
(1-\theta)^{k-1}\theta
\]

Следовательно,
\[
\xi \sim \mathrm{Geom}(\theta)
\]

\textbf{Взаимосвязи}
\begin{enumerate}
    \item Геометрическое распределение является частным случаем отрицательного биномиального распределения

Пусть случайная величина
\[
\xi \sim \mathrm{NegBin}(m,\theta)
\]
имеет отрицательное биномиальное распределение:
\[
P(\xi=x)=
\binom{x+m-1}{x}\theta^m(1-\theta)^x,
\qquad x\in\mathbb N\cup\{0\},
\]
где:
\begin{itemize}
    \item $m$ --- число успехов;
    \item $x$ --- число неудач до появления $m$-го успеха
\end{itemize}

При $m=1$ получаем:
\[
P(\xi=x)=
\binom{x}{x}\theta(1-\theta)^x
=
\theta(1-\theta)^x,
\qquad x=0,1,2,\dots
\]

Так как
\[
\binom{x}{x}=1,
\]
то
\[
P(\xi=x)=\theta(1-\theta)^x
\]

Это распределение числа неудач до первого успеха

Если же геометрическая случайная величина $\eta$ задаёт номер испытания, в котором произошёл первый успех, то
\[
\eta = \xi+1
\]

Следовательно,
\[
\eta-1 \sim \mathrm{NegBin}(1,\theta)
\]

    \item Сумма независимых геометрических случайных величин имеет отрицательное биномиальное распределение

    Пусть
    \[
    \xi_1,\dots,\xi_r
    \]
    независимы и
    \[
    \xi_i \sim \mathrm{Geom}(\theta)
    \]

    Тогда
    \[
    S=\xi_1+\dots+\xi_r
    \]
    есть номер испытания, в котором произошёл $r$-й успех

    Следовательно,
    \[
    S \sim \mathrm{NegBin}(r,\theta)
    \]
\end{enumerate}

\subsection{\texorpdfstring
{Треугольное распределение $\mathscr{L}_c=\mathrm{Triang}(0,1,\theta)$}
{Треугольное распределение Lc = Triang(0,1,theta)}}

Пусть случайная величина $\eta$ имеет треугольное распределение на $[0,1]$
с модой $\theta$, где $0<\theta<1$

Плотность:
\[
f(x)=
\begin{cases}
\dfrac{2x}{\theta}, & 0\le x\le \theta, \\[1ex]
\dfrac{2(1-x)}{1-\theta}, & \theta < x \le 1, \\[1ex]
0, & \text{иначе}
\end{cases}
\]

\textbf{Эксперимент.}
Рассматривается время выполнения лабораторной работы студентом, нормированное к интервалу $[0,1]$:
\begin{itemize}
    \item $0$ --- минимально возможное время;
    \item $1$ --- максимально возможное время;
    \item $\theta$ --- наиболее вероятное время выполнения.
\end{itemize}

Предполагается:
\begin{itemize}
    \item значения, близкие к $\theta$, встречаются чаще всего;
    \item чем дальше время от $\theta$, тем менее оно вероятно;
    \item вероятность меняется линейно
\end{itemize}

Тогда время выполнения моделируется треугольным распределением:
\[
\eta \sim \mathrm{Triang}(0,1,\theta).
\]

\textbf{Обоснование.}
Плотность:
\begin{itemize}
    \item линейно возрастает на $[0,\theta]$;
    \item достигает максимума в точке $\theta$;
    \item линейно убывает на $[\theta,1]$
\end{itemize}

Это соответствует ситуации, когда наиболее вероятным является значение $\theta$

\textbf{Взаимосвязи}
\begin{enumerate}
    \item Симметричное треугольное распределение получается как сумма двух независимых равномерных случайных величин

    Пусть
    \[
    U_1,U_2 \sim U[0,1]
    \]
    независимы

    Рассмотрим
    \[
    S=U_1+U_2.
    \]

    Тогда плотность:
    \[
    f_S(x)=
    \begin{cases}
    x, & 0\le x\le1,\\
    2-x, & 1<x\le2.
    \end{cases}
    \]

    Следовательно,
    \[
    \frac{U_1+U_2}{2}
    \sim
    \mathrm{Triang}\left(0,1,\frac12\right)
    \]

    \item Треугольное распределение возникает как свёртка равномерных распределений.

    Для независимых случайных величин плотность суммы вычисляется по формуле:
    \[
    f_S(x)=
    \int_{-\infty}^{\infty}
    f_{U_1}(t)f_{U_2}(x-t)\,dt.
    \]

    Для равномерных распределений свёртка даёт кусочно-линейную треугольную плотность

\end{enumerate}

\section{Исследование вероятностного распределения}

\subsection{\texorpdfstring
{Геометрическое распределение $\mathscr{L}_d=\mathrm{Geom}(\theta)$}
{Геометрическое распределение \mathscr{L}d = Geom(theta)}}

Пусть случайная величина $\xi$ имеет геометрическое распределение:
\[ P(\xi=k)=(1-\theta)^{k-1}\theta, \qquad k=1,2,\dots, \] где $0<\theta<1$

\subsubsection*{1. Функция распределения}

По определению:
\[F(x)=P(\xi\le x)\]

Так как $\xi$ принимает только натуральные значения, получаем:
\[F(x)=\begin{cases}0, & x<1,\\[1ex]\displaystyle\sum_{k=1}^{\lfloor x \rfloor}(1-\theta)^{k-1}\theta,& x\ge1\end{cases}\]

Сумма является суммой геометрической прогрессии:
\[\sum_{k=1}^{n}(1-\theta)^{k-1}\theta=\theta\frac{1-(1-\theta)^n}{1-(1-\theta)}=1-(1-\theta)^n\]

Следовательно,
\[F(x)=\begin{cases}0, & x<1,\\[1ex]1-(1-\theta)^{\lfloor x \rfloor},& x\ge1\end{cases}\]

\subsubsection*{2. Проверка свойств функции распределения}

\paragraph{Непрерывность справа}

Пусть $x_0\in\mathbb R$

Для дискретного распределения функция распределения имеет скачки только в целых точках и непрерывна справа:
\[\lim_{x\to x_0+0}F(x)=F(x_0)\]

Следовательно, функция распределения геометрического распределения непрерывна справа.

\paragraph{Предел при $x\to-\infty$}

Если $x<1$, то
\[F(x)=0\]

Поэтому:
\[\lim_{x\to-\infty}F(x)=0\]

\paragraph{Предел при $x\to+\infty$}

Имеем:
\[F(x)=1-(1-\theta)^{\lfloor x\rfloor}\]

Так как\[0<1-\theta<1,\] то
\[(1-\theta)^{\lfloor x\rfloor}\to0 \quad \text{при } x\to+\infty\]

Следовательно,
\[\lim_{x\to+\infty}F(x)=1\]

\subsubsection*{3. Квантили}
Квантиль уровня $\gamma$ определяется условием:
\[F(x_\gamma)\ge\gamma\]

Для геометрического распределения:
\[1-(1-\theta)^k\ge\gamma\]

Тогда:
\[(1-\theta)^k\le1-\gamma\]
\[k\ln(1-\theta)\le\ln(1-\gamma).\]

Так как
\[
\ln(1-\theta)<0,
\]
то:
\[k\ge\frac{\ln(1-\gamma)}{\ln(1-\theta)}\]

Следовательно,
\[x_\gamma=\left\lceil\frac{\ln(1-\gamma)}{\ln(1-\theta)}\right\rceil\]

Для заданных уровней:
\[x_{0.1}=\left\lceil\frac{\ln 0.9}{\ln(1-\theta)}\right\rceil\]

\[x_{0.5}=\left\lceil\frac{\ln 0.5}{\ln(1-\theta)}\right\rceil\]

\[x_{0.9}=\left\lceil\frac{\ln 0.1}{\ln(1-\theta)}\right\rceil\]

\subsection{\texorpdfstring
{Треугольное распределение $\mathscr{L}_c=\mathrm{Triang}(0,1,\theta)$}
{Треугольное распределение Lc = Triang(0,1,theta)}}

Пусть случайная величина $\eta$ имеет плотность:
\[
f(x)=
\begin{cases}
\dfrac{2x}{\theta}, & 0\le x\le\theta,\\[1ex]
\dfrac{2(1-x)}{1-\theta}, & \theta<x\le1,\\[1ex]
0, & \text{иначе}
\end{cases}
\]
где $0<\theta<1$.

\subsubsection*{1. Функция распределения}

По определению:
\[F(x)=\int_{-\infty}^{x}f(t)\,dt\]

\paragraph{При $x<0$}
\[F(x)=0\]

\paragraph{При $0\le x\le\theta$}
\[F(x)=\int_0^x \frac{2t}{\theta}\,dt=\frac{x^2}{\theta}\]

\paragraph{При $\theta<x\le1$}
\[F(x)=F(\theta)+\int_\theta^x\frac{2(1-t)}{1-\theta}\,dt\]

Так как
\[F(\theta)=\theta,\]

\[F(x)=\theta+\frac{2}{1-\theta}\int_\theta^x(1-t)\,dt\]

\[F(x)=1-\frac{(1-x)^2}{1-\theta}.\]

\paragraph{При $x>1$}
\[F(x)=1\]

Итак,
\[F(x)=
\begin{cases}
0, & x<0,\\[1ex]
\dfrac{x^2}{\theta}, & 0\le x\le\theta,\\[2ex]
1-\dfrac{(1-x)^2}{1-\theta},
& \theta<x\le1,\\[2ex]
1, & x>1.
\end{cases}\]

\subsubsection*{2. Проверка свойств функции распределения}

\paragraph{Непрерывность}

Функция распределения непрерывна как интеграл от плотности

Проверим в точке $x=\theta$:
\[\lim_{x\to\theta-0}F(x)=\frac{\theta^2}{\theta}=\theta,\]
и
\[\lim_{x\to\theta+0}F(x)=1-\frac{(1-\theta)^2}{1-\theta}=\theta\]

Тогда
\[F(\theta-0)=F(\theta+0)=F(\theta)\]

\paragraph{Предел при $x\to-\infty$}

Для $x<0$:
\[F(x)=0\]

Следовательно
\[\lim_{x\to-\infty}F(x)=0\]

\paragraph{Предел при $x\to+\infty$}

Для $x>1$:
\[F(x)=1\]

Следовательно
\[\lim_{x\to+\infty}F(x)=1\]

\subsubsection*{3. Квантили}

Квантиль уровня $\gamma$ определяется уравнением:
\[F(x_\gamma)=\gamma\]

\paragraph{Если $0<\gamma\le\theta$}

Тогда:
\[\frac{x^2}{\theta}=\gamma\]

Следовательно,
\[x_\gamma=\sqrt{\theta\gamma}\]

\paragraph{Если $\theta<\gamma<1$}
Тогда:
\[1-\frac{(1-x)^2}{1-\theta}=\gamma.\]

Отсюда:
\[(1-x)^2=(1-\theta)(1-\gamma)\]

Следовательно,
\[x_\gamma=1-\sqrt{(1-\theta)(1-\gamma)}\]

Поэтому:
\[x_{0.1}=\begin{cases}\sqrt{0.1\theta}, & 0.1\le\theta,\\[1ex]1-\sqrt{0.9(1-\theta)}, & 0.1>\theta,\end{cases}\]

\[
x_{0.5}=
\begin{cases}
\sqrt{0.5\theta}, & 0.5\le\theta,\\[1ex]
1-\sqrt{0.5(1-\theta)}, & 0.5>\theta,
\end{cases}
\]

\[x_{0.9}=1-\sqrt{0.1(1-\theta)}\]

\section{Характеристики распределений}

\subsection{\texorpdfstring
{Геометрическое распределение $\mathscr{L}_d=\mathrm{Geom}(\theta)$}
{Геометрическое распределение Ld = Geom(theta)}}

Пусть
\[P(\xi=k)=\theta(1-\theta)^{k-1},\qquad k=1,2,\dots\]

Обозначим
\[q=1-\theta\]

Тогда
\[P(\xi=k)=\theta q^{k-1}\]

\subsubsection*{1. Математическое ожидание}

По определению:
\[M\xi=\sum_{k=1}^{\infty}kP(\xi=k)\]

Подставляем распределение:
\[M\xi=\sum_{k=1}^{\infty}k\theta q^{k-1}=\theta\sum_{k=1}^{\infty}kq^{k-1}\]

Рассмотрим геометрический ряд:
\[\sum_{k=0}^{\infty}x^k=\frac1{1-x},\qquad |x|<1\]

Дифференцируем:
\[\sum_{k=1}^{\infty}kx^{k-1}=\frac1{(1-x)^2}\]

Подставляем $x=q$:
\[\sum_{k=1}^{\infty}kq^{k-1}=\frac1{(1-q)^2}\]

т.к
\[1-q=\theta,\]
то:
\[\sum_{k=1}^{\infty}kq^{k-1}=\frac1{\theta^2}\]

=>
\[M\xi=\theta\cdot\frac1{\theta^2}=\frac1{\theta}\]

Итого,
\[\boxed{M\xi=\frac1{\theta}}\]

\subsubsection*{2. Дисперсия}

Сначала вычислим второй момент:
\[M\xi^2=\sum_{k=1}^{\infty}k^2\theta q^{k-1}=\theta\sum_{k=1}^{\infty}k^2q^{k-1}\]

Используем:
\[\sum_{k=1}^{\infty}kx^k=\frac{x}{(1-x)^2}\]

Дифференцируем:
\[\sum_{k=1}^{\infty}k^2x^{k-1}=\frac{(1-x)^2+2x(1-x)}{(1-x)^4}\]
\[\sum_{k=1}^{\infty}k^2x^{k-1}=\frac{1+x}{(1-x)^3}\]

Подставляем $x=q$:
\[\sum_{k=1}^{\infty}k^2q^{k-1}=\frac{1+q}{(1-q)^3}\]

т.к
\[1-q=\theta,\]
получаем:
\[M\xi^2=\theta\cdot\frac{1+q}{\theta^3}=\frac{1+q}{\theta^2}\]

=>
\[D\xi=M\xi^2-(M\xi)^2\]

Подставляем:
\[D\xi=\frac{1+q}{\theta^2}-\frac1{\theta^2}=\frac{q}{\theta^2}\]

т.к
\[q=1-\theta,\]
получаем:
\[\boxed{D\xi=\frac{1-\theta}{\theta^2}}\]

\subsubsection*{3. Третий момент}

По определению:
\[M\xi^3=\sum_{k=1}^{\infty}k^3\theta q^{k-1}\]

Используем формулу:
\[\sum_{k=1}^{\infty}k^3x^{k-1}=\frac{1+4x+x^2}{(1-x)^4}\]

Тогда:
\[M\xi^3=\theta\cdot\frac{1+4q+q^2}{(1-q)^4}\]

т.к
\[1-q=\theta,\]
получаем:
\[M\xi^3=\frac{1+4q+q^2}{\theta^3}\]

Подставляя
\[q=1-\theta,\]
получаем:
\[\boxed{M\xi^3=\frac{1+4(1-\theta)+(1-\theta)^2}{\theta^3}}\]

\subsubsection*{4. Третий факториальный момент}

По определению:
\[M[\xi(\xi-1)(\xi-2)]\]

Вычислим:
\[M[\xi(\xi-1)(\xi-2)]=\sum_{k=1}^{\infty}k(k-1)(k-2)\theta q^{k-1}\]

Используем
\[\sum_{k=0}^{\infty}x^k=\frac1{1-x}\]

Трижды дифференцируем:
\[\sum_{k=3}^{\infty}k(k-1)(k-2)x^{k-3}=\frac{6}{(1-x)^4}\]

Умножаем на $x^2$:
\[\sum_{k=3}^{\infty}k(k-1)(k-2)x^{k-1}=\frac{6x^2}{(1-x)^4}\]

Подставляем $x=q$:
\[\sum_{k=3}^{\infty}k(k-1)(k-2)q^{k-1}=\frac{6q^2}{(1-q)^4}\]

=>
\[M[\xi(\xi-1)(\xi-2)]=\theta\cdot\frac{6q^2}{\theta^4}=\frac{6q^2}{\theta^3}\]

Итого
\[\boxed{M[\xi(\xi-1)(\xi-2)]=\frac{6(1-\theta)^2}{\theta^3}}\]

\subsubsection*{5. Производящая функция}

По определению:
\[G_\xi(s)=Ms^\xi\]

Тогда:
\[G_\xi(s)=\sum_{k=1}^{\infty}s^k\theta q^{k-1}\]

\[G_\xi(s)=\theta s\sum_{k=0}^{\infty}(qs)^k\]

Так как
\[\sum_{k=0}^{\infty}(qs)^k=\frac1{1-qs}\]
получаем:
\[\boxed{G_\xi(s)=\frac{\theta s}{1-qs}}\]

\paragraph{Первые три факториальных момента}

Дифференцируем:
\[G_\xi'(s)=\frac{\theta}{(1-qs)^2}\]

\[G_\xi'(1)=\frac{\theta}{(1-q)^2}=\frac1{\theta}\]

\[M\xi=\frac1{\theta}\]

\[G_\xi''(s)=\frac{2\theta q}{(1-qs)^3}\]

\[G_\xi''(1)=\frac{2q}{\theta^2}.\]

\[M[\xi(\xi-1)]=\frac{2(1-\theta)}{\theta^2}\]

\[G_\xi'''(s)=\frac{6\theta q^2}{(1-qs)^4}\]

\[G_\xi'''(1)=\frac{6q^2}{\theta^3}\]

\[M[\xi(\xi-1)(\xi-2)]=\frac{6(1-\theta)^2}{\theta^3}\]

\paragraph{Второй центральный момент}

Используем производящую функцию:
\[G_\xi(s)=\frac{\theta s}{1-qs},\qquad q=1-\theta\]

Ранее были найдены производные:
\[G_\xi'(s)=\frac{\theta}{(1-qs)^2},\]
\[G_\xi''(s)=\frac{2\theta q}{(1-qs)^3}\]

Подставляем $s=1$:
\[G_\xi'(1)=\frac1{\theta},\]
\[G_\xi''(1)=\frac{2q}{\theta^2}\]

т.к
\[G_\xi'(1)=M\xi,\]
\[G_\xi''(1)=M[\xi(\xi-1)],\]
то:
\[M\xi=\frac1{\theta},\]
\[M[\xi(\xi-1)]=\frac{2q}{\theta^2}\]

Выразим второй момент:
\[M\xi^2=M[\xi(\xi-1)]+M\xi\]

=>
\[M\xi^2=\frac{2q}{\theta^2}+\frac1{\theta}\]

Приводим к общему знаменателю:
\[M\xi^2=\frac{2q+\theta}{\theta^2}\]

т.к
\[q=1-\theta,\]
то
\[2q+\theta=2(1-\theta)+\theta=2-\theta\]

=>
\[M\xi^2=\frac{2-\theta}{\theta^2}\]

второй центральный момент:
\[\mu_2=M(\xi-M\xi)^2=M\xi^2-(M\xi)^2\]

\[\mu_2=\frac{2-\theta}{\theta^2}-\frac1{\theta^2}\]

Получаем:
\[\mu_2=\frac{1-\theta}{\theta^2}\]

=>
\[\boxed{\mu_2=\frac{1-\theta}{\theta^2}}\]

\paragraph{Вероятности}

Так как носитель начинается с $1$:
\[\boxed{P(\xi=0)=0}\]

\[P(\xi=3)=\theta q^2\]

=>
\[\boxed{P(\xi=3)=\theta(1-\theta)^2}\]

\[P(\xi\ge3)=\sum_{k=3}^{\infty}\theta q^{k-1}\]

\[P(\xi\ge3)=\theta q^2\sum_{k=0}^{\infty}q^k\]

т.к
\[\sum_{k=0}^{\infty}q^k=\frac1{1-q}=\frac1{\theta}\]
получаем:
\[P(\xi\ge3)=q^2\]

=>
\[\boxed{P(\xi\ge3)=(1-\theta)^2}\]

\subsubsection*{6. Характеристическая функция}

По определению:
\[\varphi_\xi(t)=Me^{it\xi}\]

Подставляем распределение:
\[\varphi_\xi(t)=
\sum_{k=1}^{\infty}
e^{itk}\theta q^{k-1},
\qquad q=1-\theta
\]

Выносим множители:
\[
\varphi_\xi(t)=\theta e^{it}
\sum_{k=0}^{\infty}(qe^{it})^k.
\]
\[\sum_{k=0}^{\infty}(qe^{it})^k=\frac1{1-qe^{it}}.\]

тогда
\[\boxed{\varphi_\xi(t)=\frac{\theta e^{it}}{1-qe^{it}}}\]

Моменты вычисляются по формуле:
\[M\xi^n=\frac1{i^n}\varphi_\xi^{(n)}(0)\]

\paragraph{Первая производная}

\[\varphi_\xi(t)=\theta e^{it}(1-qe^{it})^{-1}\]

\[\varphi_\xi'(t)=\theta\left(
ie^{it}(1-qe^{it})^{-1}
+
e^{it}\cdot
\frac{iqe^{it}}{(1-qe^{it})^2}
\right)\]

\[\varphi_\xi'(t)=
i\theta e^{it}\left(\frac1{1-qe^{it}}
+
\frac{qe^{it}}{(1-qe^{it})^2}
\right)\]

\[\varphi_\xi'(t)=i\theta e^{it}\frac{
1-qe^{it}+qe^{it}
}{
(1-qe^{it})^2
}\]

=>
\[\varphi_\xi'(t)=\frac{i\theta e^{it}}{(1-qe^{it})^2}\]

Подставляем $t=0$:
\[\varphi_\xi'(0)=\frac{i\theta}{(1-q)^2}\]

т.к
\[1-q=\theta,\]
получаем:
\[\varphi_\xi'(0)=\frac{i}{\theta}\]

=>
\[M\xi=\frac1{i}\varphi_\xi'(0)=\frac1{\theta}\]

Итого
\[\boxed{M\xi=\frac1{\theta}}\]

\paragraph{Вторая производная}

\[\varphi_\xi'(t)=i\theta e^{it}(1-qe^{it})^{-2}\]
\[\varphi_\xi''(t)=i\theta\left(
ie^{it}(1-qe^{it})^{-2}
+
e^{it}\cdot
\frac{2iqe^{it}}{(1-qe^{it})^3}
\right)\]

т.к
\[i^2=-1,\]
имеем:
\[\varphi_\xi''(t)=
-\theta e^{it}
\left(
\frac1{(1-qe^{it})^2}
+
\frac{2qe^{it}}{(1-qe^{it})^3}
\right)\]

\[\varphi_\xi''(t)
=
-\theta e^{it}\frac{
1-qe^{it}+2qe^{it}
}{
(1-qe^{it})^3
}\]

=>
\[\varphi_\xi''(t)=
-\theta e^{it}\frac{1+qe^{it}}{
(1-qe^{it})^3}\]

Подставляем $t=0$:
\[\varphi_\xi''(0)=
-\theta\frac{1+q}{(1-q)^3}.\]

т.к
\[1-q=\theta,\]
то
\[\varphi_\xi''(0)=-\frac{1+q}{\theta^2}\]

=>
\[M\xi^2=\frac1{i^2}\varphi_\xi''(0)=-\varphi_\xi''(0)\]

Тогда:
\[M\xi^2=\frac{1+q}{\theta^2}\]

т.к
\[q=1-\theta\]
получаем:
\[\boxed{M\xi^2=\frac{2-\theta}{\theta^2}}\]

\paragraph{Третья производная}

\[\varphi_\xi''(t)
=
-\theta e^{it}(1+qe^{it})(1-qe^{it})^{-3}\]

Используем правило произведения:
\[\varphi_\xi'''(t)
=
-\theta\Bigl(ie^{it}(1+qe^{it})(1-qe^{it})^{-3}\]
\[+e^{it}(iqe^{it})(1-qe^{it})^{-3}\]
\[+e^{it}(1+qe^{it})\cdot3iqe^{it}(1-qe^{it})^{-4}\Bigr)\]

\[\varphi_\xi'''(t)=-i\theta e^{it}\left[\frac{1+2qe^{it}}{(1-qe^{it})^3}+
\frac{3qe^{it}(1+qe^{it})}{(1-qe^{it})^4}\right]\]

\[\varphi_\xi'''(t)
=
-i\theta e^{it}
\frac{
(1+2qe^{it})(1-qe^{it})
+
3qe^{it}(1+qe^{it})
}{(1-qe^{it})^4}\]

\[(1+2qe^{it})(1-qe^{it})=1+qe^{it}-2q^2e^{2it}\]

\[3qe^{it}(1+qe^{it})=3qe^{it}+3q^2e^{2it}\]

Складываем:
\[1+4qe^{it}+q^2e^{2it}.\]
=>
\[\varphi_\xi'''(t)=
-i\theta e^{it}
\frac{
1+4qe^{it}+q^2e^{2it}
}{(1-qe^{it})^4}\]

Подставляем $t=0$:
\[
\varphi_\xi'''(0)
=
-i\theta\frac{
1+4q+q^2}{(1-q)^4}\]

т.к
\[1-q=\theta\]
получаем:
\[\varphi_\xi'''(0)=
-i
\frac{
1+4q+q^2
}{\theta^3}\]

=>
\[M\xi^3=\frac1{i^3}\varphi_\xi'''(0).\]

т.к
\[i^3=-i,\qquad\frac1{i^3}=i,\]

\[M\xi^3=i\left(-i\frac{1+4q+q^2}
{\theta^3}\right)\]

=>
\[M\xi^3=\frac{1+4q+q^2}
{\theta^3}\]

Подставляем
\[q=1-\theta:\]

\[\boxed{M\xi^3=\frac{1+4(1-\theta)+(1-\theta)^2}
{\theta^3}}\]

\subsection{\texorpdfstring
{Треугольное распределение $\mathscr{L}_c=\mathrm{Triang}(0,1,\theta)$}
{Треугольное распределение Lc = Triang(0,1,theta)}}

Плотность распределения:
\[
f(x)=
\begin{cases}
\dfrac{2x}{\theta}, & 0\le x\le\theta,\\[1ex]
\dfrac{2(1-x)}{1-\theta}, & \theta<x\le1,\\[1ex]
0, & \text{иначе},
\end{cases}
\qquad 0<\theta<1.
\]

\subsubsection*{1. Математическое ожидание}

По определению:
\[M\eta=\int_{-\infty}^{+\infty}xf(x)\,dx\]

Подставляем плотность:
\[M\eta=\int_0^\thetax\frac{2x}{\theta}\,dx
+
\int_\theta^1x\frac{2(1-x)}{1-\theta}\,dx\]

Первый интеграл:
\[\int_0^\thetax\frac{2x}{\theta}\,dx
=
\frac2\theta\int_0^\theta x^2\,dx\]

\[\frac2\theta
\left[
\frac{x^3}{3}
\right]_0^\theta=
\frac2\theta\cdot\frac{\theta^3}{3}=\frac{2\theta^2}{3}.\]

Второй интеграл:
\[\int_\theta^1x\frac{2(1-x)}{1-\theta}\,dx
=
\frac2{1-\theta}\int_\theta^1(x-x^2)\,dx\]

Вычисляем:
\[\frac2{1-\theta}\left[
\frac{x^2}{2}-\frac{x^3}{3}
\right]_\theta^1\]

Подставляем пределы:
\[\frac2{1-\theta}
\left(\frac12-\frac13-\frac{\theta^2}{2}+\frac{\theta^3}{3}
\right)\]

\[\frac2{1-\theta}\left(
\frac16-\frac{\theta^2}{2}+\frac{\theta^3}{3}
\right)\]

\[=\frac2{1-\theta}\cdot\frac{1-3\theta^2+2\theta^3}{6}\]

т.к
\[1-3\theta^2+2\theta^3=(1-\theta)^2(1+2\theta).\]

=>
\[=\frac{(1-\theta)(1+2\theta)}{3}.\]

Теперь складываем:
\[M\eta=\frac{2\theta^2}{3}+\frac{(1-\theta)(1+2\theta)}{3}.\]

\[(1-\theta)(1+2\theta)=1+\theta-2\theta^2\]

=>
\[M\eta=\frac{2\theta^2+1+\theta-2\theta^2}{3}\]

Получаем:
\[\boxed{M\eta=\frac{1+\theta}{3}}\]

\subsubsection*{2. Дисперсия}

Сначала вычислим второй момент:
\[M\eta^2=\int_{-\infty}^{+\infty}x^2f(x)\,dx\]

Подставляем плотность:
\[M\eta^2=
\int_0^\theta
x^2\frac{2x}{\theta}\,dx+
\int_\theta^1
x^2\frac{2(1-x)}{1-\theta}\,dx\]

Первый интеграл:
\[
\frac2\theta\int_0^\theta x^3\,dx
=
\frac2\theta\left[\frac{x^4}{4}\right]_0^\theta\]

\[=\frac2\theta\cdot\frac{\theta^4}{4}
=\frac{\theta^3}{2}.\]

Второй интеграл:
\[\frac2{1-\theta}\int_\theta^1(x^2-x^3)\,dx\]

\[\frac2{1-\theta}\left[
\frac{x^3}{3}-\frac{x^4}{4}\right]_\theta^1\]

Подставляем:
\[\frac2{1-\theta}\left(
\frac13-\frac14-\frac{\theta^3}{3}+\frac{\theta^4}{4}\right)\]

\[\frac2{1-\theta}\left(
\frac1{12}-\frac{\theta^3}{3}+\frac{\theta^4}{4}\right)\]

\[=\frac2{1-\theta}\cdot
\frac{1-4\theta^3+3\theta^4}{12}\]

т.к
\[
1-4\theta^3+3\theta^4
=(1-\theta)^2(1+2\theta+3\theta^2).\]

=>
\[=\frac{(1-\theta)(1+2\theta+3\theta^2)}{6}.\]

Тогда:
\[M\eta^2
=
\frac{\theta^3}{2}+\frac{(1-\theta)(1+2\theta+3\theta^2)}{6}.\]

Раскрываем:
\[(1-\theta)(1+2\theta+3\theta^2)=1+\theta+\theta^2-3\theta^3.\]

\[M\eta^2=\frac{3\theta^3+1+\theta+\theta^2-3\theta^3}{6}.\]

\[\boxed{M\eta^2=\frac{1+\theta+\theta^2}{6}}\]

Теперь вычисляем дисперсию:
\[D\eta=M\eta^2-(M\eta)^2\]

\[D\eta=\frac{1+\theta+\theta^2}{6}
-
\left(\frac{1+\theta}{3}\right)^2\]

\[D\eta=\frac{
3(1+\theta+\theta^2)-2(1+2\theta+\theta^2)
}{18}\]

\[=\frac{3+3\theta+3\theta^2-2-4\theta-2\theta^2
}{18}\]

\[\boxed{D\eta=\frac{1-\theta+\theta^2}{18}}\]

\subsubsection*{3. Третий момент}

По определению:
\[M\eta^3=\int_{-\infty}^{+\infty}x^3f(x)\,dx\]

Подставляем плотность:
\[M\eta^3
=\int_0^\theta
x^3\frac{2x}{\theta}\,dx
+
\int_\theta^1
x^3\frac{2(1-x)}{1-\theta}\,dx\]

Первый интеграл:
\[\frac2\theta\int_0^\theta x^4\,dx
=
\frac2\theta\left[
\frac{x^5}{5}\right]_0^\theta\]

\[=\frac{2\theta^4}{5}.\]

Второй интеграл:
\[\frac2{1-\theta}\int_\theta^1(x^3-x^4)\,dx\]

\[\frac2{1-\theta}\left[
\frac{x^4}{4}-\frac{x^5}{5}\right]_\theta^1\]

\[\frac2{1-\theta}
\left(\frac14-\frac15-\frac{\theta^4}{4}+\frac{\theta^5}{5}
\right)\]

\[\frac2{1-\theta}\left(
\frac1{20}-\frac{\theta^4}{4}+\frac{\theta^5}{5}
\right)\]

\[=
\frac2{1-\theta}\cdot
\frac{1-5\theta^4+4\theta^5}{20}\]

т.к
\[1-5\theta^4+4\theta^5=(1-\theta)^2(1+2\theta+3\theta^2+4\theta^3)\]

то
\[=\frac{(1-\theta)(1+2\theta+3\theta^2+4\theta^3)}{10}\]

\[M\eta^3=\frac{2\theta^4}{5}
+
\frac{(1-\theta)(1+2\theta+3\theta^2+4\theta^3)}{10}\]

\[(1-\theta)(1+2\theta+3\theta^2+4\theta^3)=1+\theta+\theta^2+\theta^3-4\theta^4\]

=>
\[M\eta^3=
\frac{
4\theta^4+1+\theta+\theta^2+\theta^3-4\theta^4
}{10}\]

\[
\boxed{M\eta^3
=
\frac{1+\theta+\theta^2+\theta^3}{10}}\]

\subsubsection*{4. Третий факториальный момент}

По определению:
\[M[\eta(\eta-1)(\eta-2)]\]

Раскрываем скобки:
\[\eta(\eta-1)(\eta-2)=\eta^3-3\eta^2+2\eta\]

\[M[\eta(\eta-1)(\eta-2)]=M\eta^3-3M\eta^2+2M\eta\]

Подставляем ранее найденные значения:
\[M\eta=\frac{1+\theta}{3},\]
\[M\eta^2=\frac{1+\theta+\theta^2}{6},\]
\[M\eta^3=\frac{1+\theta+\theta^2+\theta^3}{10}\]

Получаем:
\[M[\eta(\eta-1)(\eta-2)]
=
\frac{1+\theta+\theta^2+\theta^3}{10}
-
\frac{1+\theta+\theta^2}{2}
+
\frac{2(1+\theta)}{3}\]

\[=
\frac{
3(1+\theta+\theta^2+\theta^3)
-15(1+\theta+\theta^2)
+20(1+\theta)
}{30}\]

\[=
\frac{
3+3\theta+3\theta^2+3\theta^3
-15-15\theta-15\theta^2
+20+20\theta
}{30}\]

\[=\frac{
8+8\theta-12\theta^2+3\theta^3
}{30}\]

\[
\boxed{M[\eta(\eta-1)(\eta-2)]=\frac{
8+8\theta-12\theta^2+3\theta^3
}{30}}\]

\subsubsection*{5. Производящая функция}

Производящая функция вероятностей определяется только для дискретных неотрицательных целочисленных случайных величин

Так как треугольное распределение является абсолютно непрерывным распределением, производящая функция вероятностей для него не определяется

Следовательно:
\[
\boxed{
\text{Производящая функция для распределения }
\mathrm{Triang}(0,1,\theta)
\text{ не существует.}
}
\]

По этой же причине факториальные моменты и вероятности
\[
P(\eta=0),
\quad
P(\eta=3),
\quad
P(\eta\ge3)
\]
не вычисляются с помощью производящей функции

При этом:
\[
P(\eta=0)=0,
\]
так как для абсолютно непрерывной случайной величины вероятность попадания в конкретную точку равна нулю

Также:
\[P(\eta=3)=0\]

И:
\[P(\eta\ge3)=0\]

Следовательно,
\[
\boxed{P(\eta=0)=0}\]

\[\boxed{P(\eta=3)=0}\]

\[\boxed{P(\eta\ge3)=0}\]

\subsubsection*{6. Характеристическая функция}

По определению:
\[\varphi_\eta(t)=
Me^{it\eta}=\int_{-\infty}^{+\infty}e^{itx}f(x)\,dx\]

Подставляем плотность:
\[\varphi_\eta(t)=
\int_0^\theta
e^{itx}\frac{2x}{\theta}\,dx
+
\int_\theta^1
e^{itx}\frac{2(1-x)}{1-\theta}\,dx\]

Тогда:
\[\varphi_\eta(t)=
\frac2\theta\int_0^\theta
xe^{itx}\,dx
+
\frac2{1-\theta}\int_\theta^1(1-x)e^{itx}\,dx\]

по частям первый интеграл

Пусть:
\[u=x,\qquaddv=e^{itx}dx\]

Тогда:
\[du=dx,\qquadv=\frac{e^{itx}}{it}\]

\[\int xe^{itx}dx=
\frac{xe^{itx}}{it}
-
\int\frac{e^{itx}}{it}dx.\]

Получаем:
\[=\frac{xe^{itx}}{it}
-
\frac{e^{itx}}{(it)^2}\]

=>
\[
\int_0^\theta xe^{itx}dx
=\left[\frac{xe^{itx}}{it}
-
\frac{e^{itx}}{(it)^2}\right]_0^\theta\]

Подставляем пределы:
\[=\frac{\theta e^{it\theta}}{it}
-
\frac{e^{it\theta}}{(it)^2}
+\frac1{(it)^2}\]

Теперь вычислим второй интеграл:
\[\int(1-x)e^{itx}dx
=
\int e^{itx}dx
-
\int xe^{itx}dx\]

\[=
\frac{e^{itx}}{it}
-
\left(\frac{xe^{itx}}{it}
-
\frac{e^{itx}}{(it)^2}\right)\]

\[=\frac{(1-x)e^{itx}}{it}
+
\frac{e^{itx}}{(it)^2}\]

\[\int_\theta^1(1-x)e^{itx}dx
=\left[
\frac{(1-x)e^{itx}}{it}
+
\frac{e^{itx}}{(it)^2}\right]_\theta^1\]

Подставляем пределы:
\[=\frac{e^{it}}{(it)^2}
-
\frac{(1-\theta)e^{it\theta}}{it}
-
\frac{e^{it\theta}}{(it)^2}\]

\[\varphi_\eta(t)
=
\frac2\theta\left(\frac{\theta e^{it\theta}}{it}
-
\frac{e^{it\theta}}{(it)^2}
+
\frac1{(it)^2}\right)\]
\[+\frac2{1-\theta}
\left(\frac{e^{it}}{(it)^2}
-
\frac{(1-\theta)e^{it\theta}}{it}
-
\frac{e^{it\theta}}{(it)^2}\right)\]

\[\boxed{\varphi_\eta(t)
=
\frac2{(it)^2}\left(
\frac{1-e^{it\theta}}{\theta}
+
\frac{e^{it}-e^{it\theta}}{1-\theta}\right)}\]

\paragraph{Вычисление моментов с помощью характеристической функции}

Моменты вычисляются по формуле:
\[M\eta^n=\frac1{i^n}\varphi_\eta^{(n)}(0)\]

Используем характеристическую функцию:
\[\varphi_\eta(t)
=
\frac2{(it)^2}\left(\frac{1-e^{it\theta}}{\theta}
+
\frac{e^{it}-e^{it\theta}}{1-\theta}\right)\]

Разложим экспоненты в ряд Тейлора в 0 
\[e^{it\theta}=
1+it\theta+\frac{(it\theta)^2}{2}+\frac{(it\theta)^3}{6}+o(t^3)\]

\[e^{it}
=
1+it+\frac{(it)^2}{2}+\frac{(it)^3}{6}+o(t^3)\]

\[1-e^{it\theta}
=
-it\theta-\frac{(it\theta)^2}{2}-\frac{(it\theta)^3}{6}
+o(t^3),\]

и
\[e^{it}-e^{it\theta}=
it(1-\theta)+\frac{(it)^2}{2}(1-\theta^2)+\frac{(it)^3}{6}(1-\theta^3)+o(t^3)\]

\[
\frac{1-e^{it\theta}}{\theta}=
-it-\frac{i^2t^2\theta}{2}-\frac{i^3t^3\theta^2}{6}+o(t^3)\]

\[\frac{e^{it}-e^{it\theta}}{1-\theta}=it
+\frac{i^2t^2(1+\theta)}{2}
+\frac{i^3t^3(1+\theta+\theta^2)}{6}+o(t^3)\]

Складываем:
\[\frac{1-e^{it\theta}}{\theta}
+
\frac{e^{it}-e^{it\theta}}{1-\theta}\]
\[=\frac{i^2t^2}{2}
+
\frac{i^3t^3(1+\theta)}{6}
+
\frac{i^4t^4(1+\theta+\theta^2)}{24}
+
\frac{i^5t^5(1+\theta+\theta^2+\theta^3)}{120}+o(t^5)\]

тогда
\[\varphi_\eta(t)
=
\frac2{(it)^2}\left(\frac{i^2t^2}{2}
+
\frac{i^3t^3(1+\theta)}{6}
+
\frac{i^4t^4(1+\theta+\theta^2)}{24}
+
\frac{i^5t^5(1+\theta+\theta^2+\theta^3)}{120}+o(t^5)\right)\]

\[\varphi_\eta(t)
=1+\frac{it(1+\theta)}{3}+
\frac{(it)^2(1+\theta+\theta^2)}{12}+
\frac{(it)^3(1+\theta+\theta^2+\theta^3)}{60}+o(t^3)\]

Теперь дифференцируем

\paragraph{Первый момент}

\[\varphi_\eta'(t)=\frac{i(1+\theta)}{3}
+
\frac{2i^2t(1+\theta+\theta^2)}{12}
+
\frac{3i^3t^2(1+\theta+\theta^2+\theta^3)}{60}+o(t^2)\]

Подставляем $t=0$:
\[\varphi_\eta'(0)
=\frac{i(1+\theta)}{3}.\]

\[M\eta
=\frac1{i}\varphi_\eta'(0)=\frac{1+\theta}{3}.\]

\[\boxed{M\eta=\frac{1+\theta}{3}}\]

\paragraph{Второй момент}

Дифференцируем ещё раз
\[\varphi_\eta''(t)=
\frac{2i^2(1+\theta+\theta^2)}{12}
+
\frac{6i^3t(1+\theta+\theta^2+\theta^3)}{60}+o(t)\]

Подставляем $t=0$:
\[\varphi_\eta''(0)
=
\frac{i^2(1+\theta+\theta^2)}{6}.\]

Так как
\[i^2=-1\]

\[\varphi_\eta''(0)=-\frac{1+\theta+\theta^2}{6}.\]

\[M\eta^2=\frac1{i^2}\varphi_\eta''(0)
=-\varphi_\eta''(0).\]

\[M\eta^2=\frac{1+\theta+\theta^2}{6}.\]

\[\boxed{M\eta^2=\frac{1+\theta+\theta^2}{6}}\]

\paragraph{Третий момент}

Дифференцируем третий раз:
\[\varphi_\eta'''(t)
=\frac{6i^3(1+\theta+\theta^2+\theta^3)}{60}+o(1)\]

Подставляем $t=0$:
\[\varphi_\eta'''(0)=\frac{i^3(1+\theta+\theta^2+\theta^3)}{10}\]
т.к
\[i^3=-i,\]
получаем:
\[\varphi_\eta'''(0)=-\frac{i(1+\theta+\theta^2+\theta^3)}{10}\]

\[M\eta^3=\frac1{i^3}\varphi_\eta'''(0)\]

т.к
\[\frac1{i^3}=i,\]

\[M\eta^3=i\left(
-\frac{i(1+\theta+\theta^2+\theta^3)}{10}
\right)\]
\[-i^2=1,\]
имеем:
\[M\eta^3=
\frac{1+\theta+\theta^2+\theta^3}{10}.\]

\[\boxed{
M\eta^3=
\frac{1+\theta+\theta^2+\theta^3}{10}}\]

\section{Свёртка распределений}

Предполагаем, что все случайные величины независимы

\subsection{\texorpdfstring
{$\xi_1+\xi_2$, где $\xi_1,\xi_2\sim \mathrm{Geom}(\theta)$}
{xi1+xi2, где xi1, xi2 ~ Geom(theta)}}

Пусть
\[P(\xi_i=k)=\theta(1-\theta)^{k-1},\qquad k=1,2,\dots\]

Рассмотрим сумму:
\[S_2=\xi_1+\xi_2.\]

Тогда:
\[P(S_2=n)=\sum_{k=1}^{n-1}
P(\xi_1=k)P(\xi_2=n-k)\]

Подставляем распределения:
\[P(S_2=n)=
\sum_{k=1}^{n-1}
\theta(1-\theta)^{k-1}
\theta(1-\theta)^{n-k-1}\]

\[=\theta^2(1-\theta)^{n-2}\sum_{k=1}^{n-1}1\]

Так как в сумме $n-1$ слагаемых:
\[P(S_2=n)
=
(n-1)\theta^2(1-\theta)^{n-2},\qquad n=2,3,\dots\]

Следовательно
\[\boxed{\xi_1+\xi_2\sim \mathrm{NegBin}(2,\theta)}\]

\subsection{\texorpdfstring
{$\eta_1+\eta_2$, где $\eta_1,\eta_2\sim \mathrm{Triang}(0,1,\theta)$}
{eta1+eta2, где eta1, eta2 ~ Triang(0,1,theta)}}

Пусть
\[S_2=\eta_1+\eta_2,\]

\[f_\eta(x)=
\begin{cases}
\dfrac{2x}{\theta}, & 0\le x\le\theta,\\[1ex]
\dfrac{2(1-x)}{1-\theta}, & \theta<x\le1,\\[1ex]
0, & \text{иначе}\end{cases}\]

Плотность суммы независимых абсолютно непрерывных случайных величин находится по формуле свёртки:
\[f_{S_2}(x)=\int_{-\infty}^{+\infty}
f_\eta(t)f_\eta(x-t)\,dt\]

\[0\le S_2\le2.\]

Следовательно,
\[f_{S_2}(x)=0,\qquad x<0 \ \text{или}\ x>2\]

Рассмотрим 
\[0\le x\le\theta.\]

Тогда:
\[0\le t\le x,\]
и одновременно
\[0\le x-t\le\theta.\]

Следовательно, обе плотности берутся из первой ветви:
\[f_\eta(t)=\frac{2t}{\theta},\qquad f_\eta(x-t)=\frac{2(x-t)}{\theta}\]

Тогда:
\[f_{S_2}(x)=\int_0^x\frac{2t}{\theta}\frac{2(x-t)}{\theta}\,dt\]

\[=\frac4{\theta^2}\int_0^xt(x-t)\,dt\]

\[=\frac4{\theta^2}\int_0^x(xt-t^2)\,dt\]

\[=\frac4{\theta^2}
\left[\frac{xt^2}{2}-\frac{t^3}{3}\right]_0^x\]

Подставляем пределы:
\[=\frac4{\theta^2}\left(
\frac{x^3}{2}-\frac{x^3}{3}\right)\]

\[=\frac4{\theta^2}\cdot\frac{x^3}{6}.\]

\[\boxed{f_{S_2}(x)
=
\frac{2x^3}{3\theta^2},
\qquad 0\le x\le\theta}\]

Аналогично рассматриваются интервалы:
\[\theta<x\le1,
\qquad
1<x\le1+\theta,
\qquad
1+\theta<x\le2\]

На каждом из этих промежутков плотность вычисляется отдельно, поскольку меняется вид функций
\[f_\eta(t), \qquad f_\eta(x-t)\]

В результате получается кусочно-полиномиальная функция третьей степени



\subsection{\texorpdfstring
{$\xi_1+\xi_2+\xi_3$, где $\xi_i\sim \mathrm{Geom}(\theta)$}
{xi1+xi2+xi3, где xi_i ~ Geom(theta)}}

Рассмотрим сумму:
\[S_3=\xi_1+\xi_2+\xi_3.\]

\[S_3=n\]
означает, что третий успех произошёл в $n$-м испытании

Тогда:
\[P(S_3=n)
=
\binom{n-1}{2}
\theta^3(1-\theta)^{n-3},
\qquad n=3,4,\dots\]

\[\boxed{
\xi_1+\xi_2+\xi_3\sim
\mathrm{NegBin}(3,\theta)}\]

\subsection{\texorpdfstring
{$\eta_1+\eta_2+\eta_3$, где $\eta_i\sim \mathrm{Triang}(0,1,\theta)$}
{eta1+eta2+eta3, где eta_i ~ Triang(0,1,theta)}}

Пусть
\[S_3=\eta_1+\eta_2+\eta_3,\]

\[
f_\eta(x)=
\begin{cases}
\dfrac{2x}{\theta}, & 0\le x\le\theta,\\[1ex]
\dfrac{2(1-x)}{1-\theta}, & \theta<x\le1,\\[1ex]
0, & \text{иначе}
\end{cases}
\]

Плотность суммы трёх независимых случайных величин находится по формуле тройной свёртки:
\[
f_{S_3}(x)=
(f_\eta*f_\eta*f_\eta)(x).\]

То есть:
\[
f_{S_3}(x)
=
\int_{-\infty}^{+\infty}
\int_{-\infty}^{+\infty}
f_\eta(t)
f_\eta(s)
f_\eta(x-t-s)
\,dt\,ds\] 

\[0\le x\le3.\]

Следовательно,
\[f_{S_3}(x)=0,\qquad x<0 \ \text{или}\ x>3\]

Рассмотрим промежуток
\[0\le x\le\theta\]

Тогда:
\[0\le t\le x,\qquad
0\le s\le x-t,\qquad
0\le x-t-s\le\theta\]

Следовательно, во всех трёх множителях используется первая ветвь плотности:
\[f_\eta(t)=\frac{2t}{\theta},\]
\[f_\eta(s)=\frac{2s}{\theta},\]
\[f_\eta(x-t-s)=\frac{2(x-t-s)}{\theta}\]

Подставляем:
\[f_{S_3}(x)
=
\int_0^x
\int_0^{x-t}
\frac{2t}{\theta}
\frac{2s}{\theta}
\frac{2(x-t-s)}{\theta}
\,ds\,dt\]

\[=
\frac8{\theta^3}
\int_0^x
\int_0^{x-t}
ts(x-t-s)
\,ds\,dt\]

Сначала вычисляем внутренний интеграл:
\[\int_0^{x-t}s(x-t-s)\,ds\]

Раскрываем скобки:
\[=\int_0^{x-t}\bigl((x-t)s-s^2\bigr)\,ds\]

\[=\left[\frac{(x-t)s^2}{2}
-\frac{s^3}{3}\right]_0^{x-t}\]

\[=\frac{(x-t)^3}{2}
-\frac{(x-t)^3}{3}\]

Получаем:
\[=\frac{(x-t)^3}{6}\]

\[f_{S_3}(x)
=\frac8{\theta^3}\int_0^xt\frac{(x-t)^3}{6}\,dt\]

Выносим постоянную:
\[=\frac4{3\theta^3}\int_0^x
t(x-t)^3\,dt\]

Раскрываем скобки
\[(x-t)^3
=
x^3-3x^2t+3xt^2-t^3\]

Тогда
\[t(x-t)^3
=
x^3t-3x^2t^2+3xt^3-t^4\]

\[f_{S_3}(x)
=\frac4{3\theta^3}\int_0^x
\left(
x^3t-3x^2t^2+3xt^3-t^4\right)\,dt\]

\[=\frac4{3\theta^3}\left[
\frac{x^3t^2}{2}
-x^2t^3
+\frac{3xt^4}{4}
-\frac{t^5}{5}\right]_0^x\]
\[=\frac4{3\theta^3}
\left(
\frac{x^5}{2}
-x^5+\frac{3x^5}{4}
-\frac{x^5}{5}\right)\]

\[\frac12-1+\frac34-\frac15=\frac1{20}\]

\[f_{S_3}(x)=\frac4{3\theta^3}\cdot\frac{x^5}{20}\]

Получаем:\[\boxed{
f_{S_3}(x)
=
\frac{x^5}{15\theta^3},
\qquad 0\le x\le\theta}\]

На остальных промежутках:
\[\theta<x\le1,\qquad
1<x\le1+\theta,\qquad
1+\theta<x\le2,\qquad
2<x\le3,\]
плотность вычисляется аналогично, но выражения становятся более громоздкими

\subsection{\texorpdfstring
{$\sum_{i=1}^{n}\xi_i$, где $\xi_i\sim \mathrm{Geom}(\theta)$}
{sum xi_i, где xi_i ~ Geom(theta)}}

Рассмотрим сумму:
\[S_n=\sum_{i=1}^{n}\xi_i\]

\[S_n=k\]
означает, что $n$-й успех произошёл в $k$-м испытании

\[P(S_n=k)=
\binom{k-1}{n-1}
\theta^n(1-\theta)^{k-n},
\qquad k=n,n+1,\dots\]

Следовательно,
\[\boxed{
\sum_{i=1}^{n}\xi_i\sim
\mathrm{NegBin}(n,\theta)}\]

\subsection{\texorpdfstring
{$\sum_{i=1}^{n}\eta_i$, где $\eta_i\sim \mathrm{Triang}(0,1,\theta)$}
{sum eta_i, где eta_i ~ Triang(0,1,theta)}}
пусть
\[S_n=\sum_{i=1}^{n}\eta_i,\]
\[f_\eta(x)=
\begin{cases}
\dfrac{2x}{\theta}, & 0\le x\le\theta,\\[1ex]
\dfrac{2(1-x)}{1-\theta}, & \theta<x\le1,\\[1ex]
0, & \text{иначе}\end{cases}\]

Плотность суммы независимых абсолютно непрерывных случайных величин находится как $n$-кратная свёртка:
\[f_{S_n}(x)=(f_\eta*\dots*f_\eta)(x)\]

Рекуррентно:
\[f_{S_n}(x)
=
\int_{-\infty}^{+\infty}
f_{S_{n-1}}(t)
f_\eta(x-t)\,dt\]

\[0\le S_n\le n.\]

\[f_{S_n}(x)=0,\qquad x<0 \ \text{или}\ x>n\]

Рассмотрим промежуток
\[
0\le x\le\theta.
\]

Тогда
\[0\le \eta_i\le x,\]
и поэтому во всех множителях используется первая ветвь плотности:
\[f_\eta(u)=\frac{2u}{\theta}.\]

В этом случае:
\[
f_{S_n}(x)=
\frac{2^n}{\theta^n}
\int\limits_{u_1+\dots+u_n=x}
u_1u_2\dots u_n
\,du_1\dots du_{n-1}\]

\[u_i\ge0,
\qquad
u_1+\dots+u_n=x\]

Используем замену:
\[u_i=xv_i,
\qquad
v_1+\dots+v_n=1\]

Тогда
\[du_1\dots du_{n-1}=x^{n-1}dv_1\dots dv_{n-1},\] и
\[u_1\dots u_n=x^n v_1\dots v_n \]

Отсюда
\[
\boxed{
f_{S_n}(x)
=
\frac{2^n x^{2n-1}}{\theta^n}
\int\limits_{v_1+\dots+v_n=1}
v_1\dots v_n\,dv_1\dots dv_{n-1}}\]

На остальных промежутках плотность вычисляется аналогично, однако выражения становятся существенно более громоздкими, поскольку появляются различные комбинации ветвей

\subsection{\texorpdfstring
{$\xi+\eta$, где $\xi\sim \mathrm{Geom}(\theta)$, $\eta\sim \mathrm{Triang}(0,1,\theta)$}
{xi+eta, где xi ~ Geom(theta), eta ~ Triang(0,1,theta)}}

\[S=\xi+\eta,\]
где
\[P(\xi=k)=\theta(1-\theta)^{k-1},\qquad k=1,2,\dots\]
и
\[\eta\sim \mathrm{Triang}(0,1,\theta)\]

Плотность треугольного распределения
\[f_\eta(y)=
\begin{cases}
\dfrac{2y}{\theta},
& 0\le y\le \theta,
\\[8pt]
\dfrac{2(1-y)}{1-\theta},
& \theta< y\le 1,
\\[8pt]0,
& \text{иначе}
\end{cases}\]

Так как \(\xi\) — дискретная случайная величина, а \(\eta\) —
абсолютно непрерывная, плотность суммы находится по формуле
свёртки:
\[f_S(x)=
\sum_{k=1}^{\infty}
P(\xi=k)\,f_\eta(x-k)\]

Подставляем распределение геометрической случайной величины:
\[f_S(x)=\sum_{k=1}^{\infty}\theta(1-\theta)^{k-1}f_\eta(x-k)\]

Так как носитель распределения \(\eta\) равен \([0,1]\),
функция \(f_\eta(x-k)\neq0\) только при
\[0\le x-k\le1,\]
то есть
\[k\le x\le k+1\]

Следовательно, для фиксированного \(x\) в сумме остаётся
не более одного ненулевого слагаемого

Положим
\[k=\lfloor x\rfloor\]
\[f_S(x)=
\theta(1-\theta)^{k-1}f_\eta(x-k),
\qquad k\le x\le k+1\]

Если
\[k\le x\le k+\theta,\]
то
\[0\le x-k\le\theta,\]
и поэтому
\[f_\eta(x-k)=\frac{2(x-k)}{\theta}.\]

\[f_S(x)
=\theta(1-\theta)^{k-1}
\frac{2(x-k)}{\theta}
=2(1-\theta)^{k-1}(x-k).\]

Если же
\[k+\theta < x\le k+1,\]
то
\[\theta < x-k\le1,\]
и
\[f_\eta(x-k)=\frac{2(1-x+k)}{1-\theta}.\]

Тогда
\[f_S(x)
=
\theta(1-\theta)^{k-1}\frac{2(1-x+k)}{1-\theta}\]

\[\boxed{f_S(x)=
\begin{cases}
2(1-\theta)^{k-1}(x-k),
&
k\le x\le k+\theta,
\\[10pt]
\dfrac{2\theta(1-\theta)^{k-1}}{1-\theta}(1-x+k),
&
k+\theta < x\le k+1,
\end{cases}}\]
где
\[k=\lfloor x\rfloor,\qquad k=1,2,\dots\]

Следовательно, распределение суммы \(S=\xi+\eta\)
является кусочно-линейным абсолютно непрерывным распределением

\end{document}